\documentclass[11pt,a4paper]{book}
\usepackage{Qsim}



\newcommand{\myTIPTitle}{\includegraphics[width=8cm]{logo}}
\newcommand{\myTIPSubtitle}{Simulador de Arquitecturas Q}

\newcommand{\myTIPFirstAuthor}{\susi}
\newcommand{\myTIPSecondAuthor}{\tati}

\newcommand{\myTIPSubject}{Construcción de un simulador para la materia \orga}
\newcommand{\myTIPDirector}{María Nieves Dalponte}

%% 
%% Para utilizar, definir los comandos
%%  \myTIPTitle, \myTIPSubtitle
%%  \myTIPFirstAuthor\ y \myTIPSecondAuthor
%%  \myTIPSubject (el tema que se trata)
%%  \myTIPDirector
%%
\RequirePackage{graphicx}
\RequirePackage{hyperref}

\hypersetup{
pdftitle={\myTIPTitle. \myTIPSubtitle},
pdfauthor={\textcopyright\ \myTIPFirstAuthor\ y \myTIPSecondAuthor},
pdfsubject={\myTIPSubject}
}
\newcommand{\UNQlogo}[1][]{%
   \includegraphics[#1]{UNQ-logo.jpg}%
}

\newcommand{\makeCaratulaTIP}
{
                                                  \setcounter{page}{1}
%%%%%%%%%%%%%%%%%%%%%%%%%%%%%%%%%%
% The first page
%%%%%%%%%%%%%%%%%%%%%%%%%%%%%%%%%%
                                                  \phantomsection\label{portada}
                                                  \thispagestyle{empty}
                                                  \parindent=0cm 
                                                  \par\vspace{-2cm}
          \UNQlogo[scale=0.3]
                                                \begin{center}
                                                  {\large
                                                   \par\vspace{1cm}
      \textbf{Departamento de Ciencia y Tecnolog\'ia}\\ 
      \textbf{Tecnicatura Universitaria en Programaci\'on Inform\'atica}\\
												  }   
                                                \end{center}
                                                \begin{center}
                                                  \par\vspace{2cm}
                                                  \noindent%
                                                  
                                                  {\Huge\bf 
             \myTIPTitle 
                                                  \vskip 0.5cm\Large 
            \myTIPSubtitle
                                                  }{\par\vspace{1cm}
                                                  \Large\textit{ 
           \myTIPFirstAuthor                       }\par\textit{
                                                  \vskip 0.2cm\Large 
           \myTIPSecondAuthor                      }
                                                  }
                                                \end{center}
                                                  \par\vspace{2cm}
                                                \begin{center}
                                                  {\large\noindent%
                                                  \begin{tabular}{l@{\hspace{2cm}}}
      \textbf{Trabajo de Inserci\'on Profesional} \\                                                  
      \textbf{Director:} \myTIPDirector             
                                                  \end{tabular}
                                                  }\vfill\noindent%
                                                \end{center}%
}
%\author{\susi \and \tati\\ Directora: María Nieves Dalponte}
%\title{Universidad Nacional de Quilmes \\ Departamento de Ciencia y Tecnología\\\tpi\\
%\\\textbf{Simulador de Arquitecturas Q}\\Trabajo de inserción profesional}

\pagestyle{fancy}
\fancyhead{} % limpio la cabecera\nouppercase{\rightmark}


%\renewcommand{\sectionmark}[1]{\markright{#1}{}}
\fancyhead[L]{\nouppercase{\leftmark}}

\fancyhead[R]{ \thepage}
%\fancyhead[L]{\nouppercase{\leftmark \sectiontitle}}
\fancyfoot[C]{\includegraphics[width=2cm]{logo}}
\fancyfoot[R]{}

\begin{document}
\makeCaratulaTIP

%\newpage
%\chapter{Introducción}

\begin{abstract}


Una de las primeras asignaturas que debe recorrer un estudiante de la \tpi\ es \textbf{\orga}. En esta materia los estudiantes descubren los componentes funcionales que conforman un sistema de cómputos, con el fin de comprender un modelo de ejecución de programas que está presente hoy en día en la mayoría de las computadoras personales.\\

Este trabajo es el desarrollo de una herramienta que permite simular la ejecución de programas en una arquitectura teórica desarrollada por el equipo docente de la materia.\\

En este documento se presenta primero información contextual que incluye conceptos importantes y el enfoque de la materia. Luego se detalla la funcionalidad del simulador y como es el diseño orientado a objetos, siguiendo por un apartado donde se realiza un análisis del desarrollo, incluyendo el diseño de los casos de prueba. Finalmente se incluyen anexos con detalles de la arquitectura \Q, y los errores comunes en el uso del simulador.

\end{abstract}
\newpage
\section*{Agradecimientos}


\textit{''Quisiera agradecer al equipo de la materia Organización y Arquitectura por proponernos realizar su Simulador, también al equipo de Arena (Uqbar) por la gran ayuda que nos dieron con la interfaz. A Mara nuestra directora de TIP por su esfuerzo, fue muy divertido trabajar con ella. A mi compañera Tatiana por su gran gran empeño, compañerismo y ánimos, juntas pudimos lograr lo que esperábamos. A mis Padres que me financiaron la universidad para que yo pueda cursa la carrera. Por ultimo quiero decir Gracias a todas las personas que nos apoyaron llegar a nuestro objetivo. ''}\\

Susana Rosito\\

\textit{''Me gustaría agradecer en primer lugar a Mara Dalponte por su guía en este trabajo, su invaluable ayuda y buena predisposición en todo momento.
A mi compañera de equipo, Susana Rosito por su compañerismo y por trabajar codo a codo, siempre con ganas junto a mí.
Al equipo de Orga por darnos la oportunidad de realizar el simulador y facilitarnos el material de la materia para su realización.
Al equipo de Arena por facilitarnos las funcionalidades pedidas para poder realizar la interfaz.
Al director, a los docentes y a los ayudantes de la carrera por la calidad humana y la predisposición a responder consultas en todo momento, ya sea cursando o fuera de la cursada.
A la Universidad Nacional de Quilmes por apoyar a la carrera en su crecimiento continuo.
Y por último, un gran agradecimiento a mi familia, a mis amigos y a mis compañeros por apoyarme en infinitas formas desde que comencé la carrera hasta hoy.''}\\

Tatiana Molinari
%\newpage
\tableofcontents

\newpage
\parte{Contexto}
%%%%%%%%%%%%%%%%%%%%%%%%%%%%%%%%%%%%%%%%%%%%%%%%%%%%%%%%%%%
\section{Sobre la materia \orga}


Si bien los conceptos fundamentales de la ejecución de programas son independientes de las arquitecturas de computadoras comerciales, es conveniente explicar los mismos dando un marco específico. Por otro lado, las diferentes arquitecturas que subyacen los numerosos modelos disponibles en el mercado, incluyen un basto conjunto de herramientas y recursos de lenguaje para el control del funcionamiento de una computadora, pero que agregan complejidad innecesaria a la comprensión de los conceptos funcionales y la didáctica de la materia.

La propuesta de la asignatura \orga\ es utilizar la arquitectura \Q, una arquitectura \textit{assembly-like} teórica (esto es, que no existe una computadora real que la implemente) basada en el modelo de ejecución de Von Neumann y cuya característica principal es la de ser minimalista y presentarse en 'capas'. Este enfoque permite ir incorporando los conceptos de manera gradual a partir de versiones escalonadas de la arquitectura, denominadas $Q_1$, $Q_2$, $Q_3$, $Q_4$, $Q_5$ y  $Q_6$.\\


Actualmente, la arquitectura se presenta a los estudiantes mediante especificaciones formales que deben analizar y comprender para utilizar el lenguaje de programación y resolver los problemas que se les plantean en la práctica. Entendemos que les resulta de gran utilidad incorporar una herramienta que les permita probar sus ejercicios de una manera automatizada y es por eso que se desarrolló un simulador para esta arquitectura.

%%%%%%%%%%%%%%%%%%%%%%%%%%%%%%%%%%%%%%%%%%%%%%%%%%%%%%%%%%%
\section{Conceptos importantes}

\subsection{Enfoque de Von Neumann}

El matemático John Von Neumann en el año 1945 se encontraba colaborando en el proyecto ENIAC (\textit{Electronic Numerical Integrator And Computer}, primer computadora electrónica de propósito general, diseñada para ser utilizada por el ejército norteamericano. La ENIAC podía ser programada para realizar operaciones complejas e incluso decisiones, interacciones y subrutinas, pero la tarea de resolver un problema y volcarlo en la máquina era tan complejo que podía tomar semanas. Luego que el programa era diseñado en papel, el proceso de representarlo en la máquina ENIAC mediante la manipulación de cables e interruptores tomaba varios días. Entonces Von Neumann se comenzó a interesar por la problemática que significaba la la necesidad de reconfigurar la máquina para cada nueva tarea y tan sólo cuatro años más tarde propone y desarrolla una solución a este problema que se basaba en almacenar la información sobre las operaciones a realizar en la misma memoria utilizada para los datos, a partir de su codificación en código binario al igual que los datos.\\

Además, este enfoque generaliza la organización de las computadoras distinguiendo en tres partes interconectadas: La CPU (con la unidad aritmético-lógica o ALU y la unidad de control) la memoria, y un módulo de entrada/salida. La interconexión es llevada a cabo por un bus de sistema que proporciona un medio de transporte de los datos entre las distintas partes. \\

Con la propuesta de este modelo, Von Neumann incorpora el concepto de \textbf{programa almacenado} en memoria. Con esta idea, el programa se codifica de cierta manera para que pueda ser almacenado en memoria principal y posteriormente pueda ser ejecutado quizás múltiples veces. De esta manera, la lógica del programa puede ser ''recordada'' y el programa toma un valor mayor, a diferencia de lo que ocurría hasta entonces, donde el programa se reflejaba en un conjunto de configuraciones de cables aplicadas a los equipos. Esto implica una separación entre el mecanismo de ejecución (el \textit{hardware}) y la lógica de computo o instrucciones (el \textit{software}). La codificación en binario de las instrucciones de un programa se denomina \textbf{\codmaq}.\\

Por otro lado este tipo de diseño, que permite un programa almacenado, también da la posibilidad de que la ejecución de las instrucciones modifique el código máquina del mismo u otro programa. Por ejemplo un programa podría modificar o incrementar las referencias a las direcciones de memoria que tenga en algunas instrucciones y luego volver a ejecutar dichas instrucciones con el fin de procesar celdas diferentes de memoria cada vez. Esta característica es potente pero presenta un alto riesgo pues las modificaciones en los programas podía ser algo perjudicial, por accidente o por diseño.

\subsection{Organización de la computadora}

La CPU (Unidad Central de Procesamiento del inglés: Central Processing Unit), es el componente principal y el encargado ejecutar los programas y procesar los datos. La CPU contiene otros componentes de importancia tales como la \UC, algunos registros de uso específico como el contador de programa (PC o \textit{Program Counter}), el registro de instrucción (IR - \textit{Instruction Register}) y el \SP\ (SP - \textit{Stack pointer}), otros registros de uso general y la Unidad Aritmético-Lógica (ALU).\\

La \UC\ dirige el ciclo de ejecución de cada instrucción, pidiendo la lectura de celdas de memoria donde esta alojada, decodificándola  y ejecutándola luego en colaboración con los otros componentes del sistema: si es una operación lógica o aritmética le ordena a la ALU su ejecución, si es de movimiento de datos colabora con la memoria ó el módulo de Entrada/Salida.\\

Entre los registros de uso específico, los más importantes son el \textbf{\PC}, el \textbf{\IR}, el \textbf{\SP} y los \textbf{\FLAGS}. El \PC\ es un registro que indica la posición de memoria donde estará la siguiente instrucción que debe ejecutarse. Luego de completar el ciclo de ejecución de una instrucción, el PC se incrementa en función de la cantidad de celdas que ocupa el código máquina de esta. 
El registro de instrucción, \textbf{\IR}, contiene el código máquina de la instrucción actual una vez que la misma es leída de memoria para luego decodificarla y ejecutarla.
El \SP\ es un registro que indica la posición de memoria donde está el dato que representa a la última dirección de memoria que fue guardada al llamar a una sub-rutina, para que cuando esta termine, el programa pueda seguir ejecutándose a partir de ese punto.
Los \FLAGS\ representan características del último computo realizado como si el resultado fue negativo, cero, etc.

El diseño de cada arquitectura ofrece un conjunto diferente de registros de uso general para ser usados en los programas. Estos registros son elementos de memoria de alta velocidad y poca capacidad que pueden ser utilizados como variables en los programas. Es importante marcar que pueden almacenar tanto datos como direcciones de memoria.\\

La \ALU, recibe su nombre de las siglas en inglés de \textit{Arithmetic and Logic Unit}. La ALU es un circuito digital que lleva a cabo operaciones aritméticas (suma, resta, multiplicación, división) y las operaciones lógicas como la negación, disyunción, conjunción, etc, entre dos cadenas binarias que son interpretadas como números o valores lógicos.\\

La memoria es un conjunto de celdas numeradas. La numeración de cada celda la identifica inequívocamente por lo cual a esta numeración se le llama dirección. En cada celda de la memoria se pueden almacenar datos o instrucciones en forma de cadenas binarias y este contenido puede leerse y modificarse. En la memoria es donde se alojan los programas que luego serán ejecutados.\\

El bus de sistema es el encargado de transferir los datos entre los componentes de la computadora..La unidad de control al pedir un contenido de una dirección de memoria lo hace a través del bus, y similarmente mismo cuando desea escribir en memoria, y lo hace a partir de que la \UC\ pide la lectura o escritura de celdas de memoria o puertos de entrada/salida.

\subsection{Ejecución de un programa}

La función de una computadora es la ejecución de programas. Los programas se encuentran almacenados en memoria y consisten en una secuencia de instrucciones y es la \UC\ es quien se encarga de ejecutar dichas instrucciones implementando un \textbf{\ciclo}. Para ser almacenadas en memoria, las instrucciones deben codificarse en cadenas binarias (secuencias de ceros y unos) que no son legibles para las personas pero son tales que la \UC\ las puede interpretar y traducir en acciones. Por eso para saber de qué instrucción se trata, y cuáles son los valores o variables (celdas de memoria o registros) que están involucrados, la \UC\ toma el código máquina de la instrucción y verifica los códigos de operación y modos de direccionamiento. La ejecución de instrucciones se divide en tres etapas importantes: 

\begin{enumerate}
\item \BI
\item \DI
\item \EI
\end{enumerate}
\graf[7cm]{ciclo}{\ciclo}


Al principio de cada \ciclo se lleva a cabo la \BI\ durante la cual se leen las celdas que contienen el \codmaq\ de la instrucción y para lo cual se utiliza el \PC\ que mantiene la dirección de la siguiente instrucción a ejecutar.
El código máquina de la instrucción leída que está en la forma de cadena binaria se carga dentro de otro registro de la CPU, llamado registro de instrucción (IR).\\

Durante la \DI, la \UC\ determina que operación se debe llevar a cabo, y finalmente durante la \EI la \UC realiza el efecto esperado para esa operación, buscando los operandos y modificando la memoria o los registros como resultado final y el ciclo vuelve a comenzar

%%%%%%%%%%%%%%%%%%%%%%%%%%%%%%%%%%%%%%%%%%%%%%%%%%%%%%%%%%%
\section{Estado del arte}

A través de los años de la carrera \tpi, en la materia \orga\ se analizaron distintos enfoques y herramientas para desarrollar los conceptos relacionados a la ejecución de programas en una computadora.\\

Inicialmente se utilizó un simulador de código abierto para la arquitectura Intel 8085, que ofrecía una funcionalidad bastante completa, pero una interfaz que no resultaba del todo intuitiva.
Este simulador se eligió por tratarse de un lenguaje assembler mas reducido, con un repertorio de instrucciones y modos de direccionamiento mas pequeño que contaba con un entorno de prueba (el simulador propiamente dicho) para facilitar la didáctica de la programación en lenguaje ensamblador, pero posteriormente se entendió que las características de la arquitectura no eran las adecuadas para la enseñanza de los contenidos y se descartó.\\
 
Entonces el equipo docente eligió definir una arquitectura teórica que proveyera solamente lo necesario para cumplir con los objetivos de la materia e inspirados en un caso similar de la Universidad de Buenos Aires, se definió la arquitectura \textbf{QARQ}, muy similar a la que se presenta en la sección \ref{apendiceQ}\\


Posteriormente, el equipo docente propuso un cambio en la secuencia didáctica que requirió la división de la especificación de la arquitectura QARQ en varias partes, donde cada una recibe el nombre de \textbf{Arquitectura Qi} y agrega una nueva funcionalidad (instrucciones o modos de direccionamientos) a la versión anterior, construyendo una arquitectura en capas. Sin embargo todos los ejercicios de las arquitecturas Qi se siguen haciendo en papel. \\

Actualmente se busca incorporar el \textbf{Simulador Qsim} en la materia con el objetivo de que los alumnos puedan visualizar el funcionamiento de una computadora al mas mínimo detalle, a través de la ejercitación del \ciclo.

%%%%%%%%%%%%%%%%%%%%%%%%%%%%%%%%%%%%%%%%%%%%%%%%%%%%%%%%%%%
\section{Arquitecturas Q}

Las versiones de la arquitectura Q están pensadas para incorporar funcionalidades de manera que la curva de aprendizaje sea adecuada para los alumnos, siendo paulatina e incremental, es decir, cada arquitectura $Q_{i+1}$ agrega más funcionalidad (ya sean instrucciones nuevas o modos de direccionamiento) a la arquitectura  $Q_{i}$ anterior (ver figura \ref{ArquitecturasQ})

\graf[13cm]{ArquitecturasQ}{Organización en capas de la arquitectura \Q}

%%%%%%%%%%%%%%%%%%%%%%%%%%%%%%%%%%%%%%%%%%%%%%%%%%%%%%%%%%%



\part{Simulador \qsim}


El simulador desarrollado es una herramienta de utilización sencilla ya que se asume que los alumnos de esta materia están en la etapa inicial de la carrera y se pretende transmitir los conceptos sin distraerlos con detalles de uso y configuración.

\section{Funcionalidad del simulador}

La funcionalidad del simulador puede caracterizarse mediante las siguientes partes importantes:

\begin{itemize}
\item Chequeo de sintaxis de los programas escritos en el lenguaje Q
\item Ensamblado del código fuente de un programa en su correspondiente código máquina
\item Cargado en memoria del código máquina
\item Ejecución paso a paso de un programa cargado en memoria
\end{itemize}

\subsection{Chequeo de sintaxis}\label{parser}

El simulador provee al alumno de un editor de texto en el cual escribirá el programa en un lenguaje Qi, que desea cargar en memoria y ejecutar.
Una vez que el usuario haya terminado la escritura, al momento de cargar el programa, el simulador utilizará un \textit{parser}\footnote{Un parser (o analizador sintáctico) es una de las partes del compilador que transforma su entrada de texto plano en este caso a objetos del modelo.} para detectar errores de sintaxis, tales como la falta de una coma o un corchete, o la presencia de símbolos que no pertenecen al lenguaje (como por ejemplo signos de pregunta y símbolos matemáticos); o bien errores semánticos como la combinación incorrecta de elementos del lenguaje, por ejemplo: modos de direccionamiento mal ubicados.
El parser solo revisará lo escrito por el alumno y de acuerdo a las gramática del lenguaje, mostrará alguno de los siguientes estados:

\begin{description}
\item[OK] Este mensaje se obtiene cuando no hubo ningún error de sintaxis. Si se da este resultado, es posible continuar con el ensamblado y cargado en memoria.
\item[SyntaxError] Este mensaje de error se obtiene cuando en alguna línea del programa se detectó algún error de sintaxis o de semántica, como se describió arriba. Cuando ocurre este error se lo acompaña con una descripción lo mas detallada posible para que el alumno detecte donde ocurrió y pueda corregirlo. Un programa con errores no puede ser ensamblado y cargado en memoria.
\end{description}

\subsection{Ensamblado}

Una vez que el programa es sintácticamente válido es posible traducir el código fuente del programa en código máquina (representado en cadenas binarias). Para esto se respeta un formato de instrucción que indica cómo se codifica cada operación y los operandos. \\
Mas detalle al respecto de este proceso en el apéndice \ref{apendiceQ}.

\subsection{Cargado en memoria}

Una vez ensamblado, la representación binaria (o código máquina) del programa será cargado en memoria a partir de una ubicación (celda de memoria) que el alumno puede elegir. Esto permite visualizar el contenido de la memoria (con el programa cargado) y el estado de los registros de la CPU. 

Durante la carga del programa en memoria puede ocurrir que el programa no cuente con el espacio suficiente a partir de la ubicación elegida ya que ocupa más celdas que las que se encuentran disponibles, ya que como se especifica en el apéndice \ref{apendiceQ}, la memoria disponible tiene un tamaño limitado y por este motivo la alocación en memoria del código máquina puede exceder el espacio disponible a partir de la celda inicial anteriormente elegida. Si por el contrario, no se produce este error, el alumno podrá ver el programa cargado en memoria exitosamente.

\subsection{Ejecución paso a paso}

Se provee la funcionalidad de la ejecución paso a paso ya que se desea que el alumno pueda experimentar y así comprender los pasos del \ciclo. El paso a paso que provee el simulador consiste en las siguientes etapas pertenecientes al \ciclo:

\begin{enumerate}
\item \textbf{Búsqueda de instrucción:} El alumno podrá visualizar el valor que contiene PC (Program counter) donde se encuentra la dirección de la celda en memoria que contiene la próxima instrucción a ejecutar (por ejemplo, en caso de ser la primer instrucción del programa recién cargado, el pc tendrá la dirección de memoria elegida por el alumno para iniciar el cargado del programa en memoria). El simulador, toma de la memoria el código maquina correspondiente a la instrucción que comienza en esa dirección tomada de PC (una instrucción puede ocupar más de una celda de memoria) y los guarda en el \IR\ (\textit{Instruction Register}). Será observable también para el alumno el incremento del registro PC, tanto como celdas ocupe la instrucción actual, lo que conceptualmente es, preparar el contexto de ejecución para tomar la siguiente instrucción.

\item  \textbf{Decodificación:}
En la decodificación el interprete se encarga de desensamblar el código máquina (abreviado en hexadecimal) que ya fue ubicado en el \IR\ para mostrar el código fuente de la instrucción actual con sus respectivos operandos. Si el programa escrito por el alumno es sintacticamente y conceptualmente correcto, este paso le permite comprobar que la instrucción actual es la que él mismo escribió y no otra, visualizándola en pantalla. \\

La decodificación con desensamblado, es decir con la reconstrucción del código fuente, permite al alumno experimentar otros escenarios y efectos laterales, entre los cuales podemos enumerar:

\begin{itemize}
\item Si la ejecución paso a paso excede los límites del programa, pueden tomarse instrucciones de otra rutina y procesarse como una nueva instrucción.
\item Si en cambio, se intenta ejecutar el contenido de una celda con datos (y no una instrucción) podrá ocurrir que se encuentre una instrucción invalida (por ejemplo, una combinación incorrecta de modos de direccionamiento y códigos de operación) y el alumno verá el mensaje de error pertinente.
\item Que el programa sobrescriba su mismo código máquina.
\end{itemize}

Mas detalle de estas situaciones en la sección \ref{erroresconceptuales}.

\item  \textbf{Ejecución}
La ejecución lleva a cabo los efectos de la instrucción y muestra en pantalla los cambios en el estado de ejecución: contenidos de memoria, puertos, registros y flags. Esto quiere decir que es en esta etapa cuando se lleva a cabo el almacenamiento de resultados que, si es necesario, guardará el valor resultante de la operación descripta por la instrucción en el operando destino. Esto cambiará el valor de una celda de memoria o de un registro y será visto en pantalla por el alumno.
\end{enumerate}


%%%%%%%%%%%%%%%%%%%%%%%%%%%%%%%%%%%%%%%%%%%%%%%%%%%%%%%%%%%5
\section{Implementación}

%%%%%%%%%%%%%%%%%%%%%%%%%%%%%%%%%%%%%%%%%%%%%%%%%%%%%%%%%%%
\subsection{Tecnología utilizada}
En la presente sección se indica la tecnología elegida para la implementación del simulador, justificando dichas elecciones en cada caso.

\begin{itemize}


\item  \textbf{Lenguaje Scala.}
Elegimos el lenguaje Scala para realizar el simulador ya que durante las cursadas de las asignaturas de TPI no tuvimos la oportunidad de profundizar el dominio de este lenguaje ni aprovechar las ventajas que este ofrecía al combinar el manejo de objetos y las características de un lenguaje funcional. Es por ello que descartamos la elección de un lenguaje con el que estábamos mas familiarizadas, como por ejemplo Java.

\item  \textbf{Framework Arena.}
Utilizamos el framework Arena para realizar la interfaz de usuario del simulador porque es una herramienta de código abierto que tuvimos la oportunidad de conocer en la materia \ui. Al poder ser combinado con \textbf{Scala} nos pareció una buena oportunidad de explotar lo que nos ofrecía para simplificar la definición de la interfaz de usuario, permitiéndonos así enfocarnos en la implementación del modelo. 

\item  \textbf{Eclipse.}
Se eligió utilizar el entorno de programación Eclipse ya que es una herramienta multiplataforma, lo que nos permitió trabajar en diferentes sistemas operativos y con la cual estábamos familiarizadas.  Por otro lado, la comunidad provee pluggins para manejar proyectos para Scala.

\item  \textbf{Git}
Elegimos git como repositorio externo para trabajar colaborativamente durante el desarrollo, dado que es una herramienta que muy extendida en los desarrollos de software libre.

\end{itemize}

%%%%%%%%%%%%%%%%%%%%%%%%%%%%%%%%%%%%%%%%%%%%%%%%%%%%%%%%%%%
\subsection{Diseño Orientado a Objetos}

\subsubsection{ALU}
Como se observa en la figura \ref{ALU} la ALU tiene toda la responsabilidad en la ejecución de operaciones matemáticas y lógicas, además del cómputo de los flags luego de cada operación. 

\graf{ALU}{Diagrama de clase de la \ALU}
%%%%%%%%%%%%%%%%%%%%%%%%%%%%%%%%%%%%%%%%

\subsubsection{Bus de entrada y salida, memoria y puertos}
Como se observa en la figura \ref{BusEntradaSalida_Memoria_CeldasPuertos} el Bus de entrada y salida tiene la responsabilidad de derivar según donde corresponda (Memoria o Puertos) la escritura o lectura de un dato. Para ello conoce a una instancia de la clase Memoria y a otra de la clase CeldasPuertos. 
Ambas clases conocen una colección de instancias de la clase Celda, y cada Celda a su vez conoce un dato: una instancia de la clase W16 que representa al dato almacenado en una celda de memoria o en un puerto.  

\graf{BusEntradaSalida_Memoria_CeldasPuertos}{Diagrama de las clases del la BusEntradaSalida, Memoria, CeldasPuertos y Celda}
%%%%%%%%%%%%%%%%%%%%%%%%%%%%%%%%%%%%%%%%

\subsubsection{CPU}
Como se observa en la figura \ref{CPU}, la CPU conoce a la ALU, contiene los registros IR y PC, los flags (V,Z,C,N) y los ocho registros de uso general (R0...R7). La responsabilidad de la CPU es actualizar los flags, los registros, actualizar el PC y el IR, y ser la conexión con la ALU.

\graf{CPU}{Diagrama de clase de la CPU}
%%%%%%%%%%%%%%%%%%%%%%%%%%%%%%%%%%%%%%%%

\subsubsection{Intérprete} \label{claseinterprete}
Como se observa en la figura \ref{Interprete} el Interprete es un \textit{singleton} que tiene la entera responsabilidad de construir un objeto que representa una instrucción determinada a partir de la decodificación de las celdas de memoria que ocupa su código máquina. Se ocupa de la decodificación de la instrucción.

\graf{Interprete}{Diagrama de clase del Interprete}
%%%%%%%%%%%%%%%%%%%%%%%%%%%%%%%%%%%%%%%%

\subsubsection{Modos de direccionamiento y W16}

Esta jerarquía de clases permite controlar el acceso a los operandos, mediante un objeto ModoDireccionamiento que controla el acceso a un objeto W16, que es quien contiene el valor propiamente dicho. 
Como se observa en la figura \ref{ModoDireccionamiento} los modos de direccionamiento extienden del \textit{trait (un tipo abstracto, utilizado como modelo para estructurar programas orientados a objetos) ModoDireccionamiento, donde se encuentran declarados mensajes necesarios para manejar todas las subclases de manera polimórfica. Entre estos mensajes podemos enumerar:
\begin{itemize}
\item  \textbf{representacionString:}
Devuelve la representación en string como código fuente, por ejemplo la representación de un objeto
 \textbf{ADD(R0,R7)}, sería: \code{ADD R0, R7} este mensaje se utiliza en la etapa de desensamblado.
\item  \textbf{codigo:}
Retorna el string que representa al código del modo de direccionamiento, por el ejemplo, el código de modo de direccionamiento del R7 es \code{100111}. este mensaje se utiliza en la etapa de ensamblado.
\item  \textbf{getValorString:}
Retorna el dato almacenado en el operando. En el caso de un Inmediato que sea \code{FF56}, devolverá el string "FF56", y en el caso de cualquier registro, retornara el valor que represente su W16.
\end{itemize}


La clase W16 que también esta en la figura \ref{ModoDireccionamiento}, representa el dato que es guardado en memoria. Tiene la capacidad de incrementarse, decrementarse, sumar una entero, devolver su representación binaria y su valor en decimal.\\

Los modos de direccionamiento diferentes a Inmediato y Registro, conocen otro modo de direccionamiento que encapsula al objeto \textbf{W16} según corresponda, es decir:
\begin{itemize}
\item \textbf{RegistroIndirecto} conoce una instancia de Registro.
\item \textbf{Directo} conoce conoce una instancia de Inmediato.
\item \textbf{Indirecto} conoce conoce una instancia de Directo.
\end{itemize}

La clase \textbf{Etiqueta} representa las etiquetas creadas por el alumno cuando realiza el programa. Cuando el mismo es cargado en memoria, en función de cual sea la celda de inicio y cuanto ocupen las instrucciones, se calcula la dirección de memoria a la que hace referencia y luego se la descarta reemplazándola por un modo de direccionamiento Inmediato.

\graf[15cm]{ModoDireccionamiento}{Diagrama de clase de la jerarquía de los modos de direccionamiento}
%%%%%%%%%%%%%%%%%%%%%%%%%%%%%%%%%%%%%%%%

\subsubsection{Instrucciones}
Como se observa en la figura \ref{Instruccion} las Instrucciones están jerarquizadas en:

\begin{itemize}
\item Instrucciones de un operando.
\item Instrucciones de dos operandos.
\item Instrucciones de sin operandos.
\item Saltos condicionales.
\end{itemize}

Se realizó dicha jerarquía para permitir la fácil incorporación de nuevas instrucciones como subclases de la clase que corresponda ya que de ese modo se reutiliza comportamiento tal como la manera de mostrarse (en términos de código fuente) y de codificarse (en términos de \codmaq).

\graf{Instruccion}{Diagrama de clase de la Instrucción}

\begin{description}
\item[Instrucciones sin operandos] Como se observa en la figura \ref{Instruccion_SinOperandos} la única instrucción sin operandos implementada en la arquitectura Q es la instrucción \code{RET}. A pesar de esto, se eligió hacer una jerarquía para que luego se facilite la escalabilidad del modelo, permitiendo la inserción de nuevas instrucciones sin operandos.

\graf[4cm]{Instruccion_SinOperandos}{Detalle de clase Instruccion\_SinOperandos}

\item[Instrucciones de un operando] 
Como se observa en la figura \ref{Instruccion_UnOperando} y siguiendo con el criterio mencionado antes sobre la jerarquía de instrucciones, puede haber dos tipos de instrucciones de un operando:

\begin{itemize}
\item Un operando origen.
\item Un operando destino.
\end{itemize}

Ambas clases de instrucciones tienen un solo operando. La diferencia entre ellas es la lógica de ejecución y la manera en la que se ensamblan y desensamblan ya que sus formato de instrucción difiere (ver apéndice \ref{apendiceQ}).

\graf[15cm]{Instruccion_UnOperando}{Diagrama de clase de la Instruccion\_UnOperando}


\item[Instrucciones de dos operandos]

Como se observa en la figura \ref{Instruccion_DosOperandos} la jerarquía de clases de instrucciones de dos operandos es la más amplia por tener mayor cantidad de instrucciones. Todas comparten el comportamiento para la decodificación e interpretación, además de la lógica de impresión.

\graf[15cm]{Instruccion_DosOperandos}{Diagrama de clase de la Instruccion\_DosOperandos}

\end{description}



%%%%%%%%%%%%%%%%%%%%%%%%%%%%%%%%%%%%%%%%

\subsubsection{Programa}

Como se observa en la figura \ref{Programa} la clase Programa conoce un grupo de Instrucciones (las instrucciones que lo componen). Las instancias son creadas por el \textit{parser} (ver sección \ref{parser}), luego se calculan las etiquetas (si es que las tiene) y finalmente cuando es cargado en memoria la instancia de programa es desechada ya que no vuelve a usarse en ningún momento de la ejecución.

\graf{Programa}{Diagrama de clase de la Programa}
%%%%%%%%%%%%%%%%%%%%%%%%%%%%%%%%%%%%%%%%

\subsubsection{Simulador}

Esta es la clase principal del modelo y la encargada de coordinar la ejecución del programa paso a paso. Como se observa en la figura \ref{simulador} la clase \textbf{Simulador} conoce a una instancia de la clase \textbf{CPU}, a una instancia de la clase \textbf{BusEntradaSalida} y a una instancia de la clase \textbf{Instruccion}, que representa a la instrucción que se esta ejecutando en ese momento.\\

La clase Simulador tiene la responsabilidad de obtener el \codmaq\ de la siguiente instrucción, colaborando con el Interprete (ver \ref{claseinterprete}), calcular las etiquetas de un programa, cargar el programa en memoria y los datos en registros,  ejecutar las instrucciones o delegar su ejecución al objeto \textbf{ALU} que conoce a través de la \textbf{CPU} según corresponda, y guardar datos en memoria o registros (almacenamiento de resultados).

\graf[15cm]{simulador}{Diagrama de clase de la Simulador}
%%%%%%%%%%%%%%%%%%%%%%%%%%%%%%%%%%%%%%%%

%Como se observa en la figura \ref{}

%\graf{}{Diagrama de clase de la \}
%%%%%%%%%%%%%%%%%%%%%%%%%%%%%%%%%%%%%%%%

%[width=0.7\textwidth]




\parte{Evaluación del desarrollo}

%%%%%%%%%%%%%%%%%%%%%%%%%%%%%%%%%%%%%%%%%%%%%%%55
\section{Dificultades encontradas}

\subsection{Dificultades presentadas por el dominio}
Las dificultades del dominio estuvieron relacionadas a la comprensión no solamente del modelo de arquitectura Q si no también a su propósito didáctico, ya que el objetivo del simulador no es solamente la simulación de la arquitectura y la ejecución de programas, sino también el proveer a los alumnos la capacidad de ejercitar situaciones conceptualmente erróneas, por lo que se requería una comprensión didáctica del problema más allá de la especificación de las arquitecturas Q. 

\subsection{Dificultades de diseño}
En primera instancia se optó por implementar un modelo de objetos que utilizaba un objeto de la clase \textbf{Programa} a lo largo de toda la ejecución. Esto permitía evitar la lectura de la memoria principal para obtener las instrucciones a ejecutar, solicitando cada instrucción al objeto de la clase Programa, sobreviviendo así las distintas etapas una vez que fue creado por el \textit{parser}. Este enfoque evitaba la necesidad de un objeto que tuviera que interpretar el código máquina alojado en la memoria, evitaba la nueva creación de instancias de la clase \textbf{Instrucción} y simplificaba en gran medida el modelo ya que la memoria era sólo una clase que contenía datos que se reflejaban en pantalla y no se la utilizaba en el la búsqueda de instrucciones. Luego entendimos que un programa no sólo podía modificar su entorno al ser ejecutado (otras celdas de memoria que no ocupen su código maquina, celdas de puertos, registros, etc) si no que también podría sobrescribir su código maquina (ya sea con ese propósito o sólo por un error conceptual), o bien, que el alumno debía tener la posibilidad de seguir ejecutando más allá del código máquina alojado en memoria o más. A partir de eso fue necesario corregir gran parte del modelo agregando una clase denominada \textbf{Intérprete}, cuya responsabilidad es interpretar la siguiente instrucción alojada en memoria para que luego sea ejecutada, otorgando más responsabilidad a la clase \textbf{Memoria} y descartando el objeto instancia de \textbf{Programa} una vez que éste es cargado en memoria con éxito.\\

Por otro lado, tuvieron que solicitarse extensiones al equipo de desarrolladores de Arena para poder implementar la interfaz del simulador utilizando dicho framework. Entre dichas extensiones podemos enumerar:

\begin{itemize}
\item \textbf{FileSelector}: que permite la carga de archivos de código fuente.
\item \textbf{CodeEditor} (O actualmente llamado KeywordTextArea): Es un área en la ventana donde se visualiza el el código fuente de los programas   una vez que han sido elegidos del navegador de archivos.
\item \textbf{Bindings} contra el background de componentes y celdas de una tabla, para permitir la visualización de cambios de colores en la memoria a partir de que se recorren las celdas del programa o se realiza algún cambio en el contenido de una celda.
\item \textbf{TexBox multiline} es un área de texto que puede navegarse (utilizando el \textit{scroll}) para ser utilizado como consola de devolución.
\item Icono para la aplicación, solo por cuestiones estéticas.
\end{itemize}

%%%%%%%%%%%%%%%%%%%%%%%%%%%%%%%%%%%%%%%%%%%%%%%55
\section{Casos de prueba}

En esta sección se describen los casos de prueba de la aplicación realizados para ilustrar el alcance de la aplicación y su robustez.

\subsection{Chequeo de sintaxis en la distintas Qi}

El chequeo de sintaxis de cada versión de la \textbf{Arquitectura Q} se lleva a cabo procesando algunos programas con el \textit{parser}, y con este objetivo se confeccionaron  dos casos de prueba por cada Qi:
\begin{enumerate}
\item \textbf{Chequear programa Qi válido}

Se utiliza como parámetro para el objeto de la clase \textbf{Parser} un programa Qi válido, es decir que es sintáticamente correcto en la arquitectura Qi, y éste comprobará la sintaxis del programa acorde a la arquitectura Qi seleccionada. Por ejemplo, un programa sintácticamente válido en Q6 puede no ser sintácticamente válido en Q1 si utiliza modos de direccionamiento o instrucciones que Q1 no incluye. Al terminar de chequear el programa, el objeto Parser retorna un objeto instancia de la clase Programa con la lista de instrucciones que fueron ensambladas. 
Se toma dicho resultado y como ultimo paso se compara con el programa esperado (que concuerda en objetos con el programa en código fuente escrito en el input).


\item \textbf{Chequear programa Qi es invalido} 

Se utiliza como parámetro para el objeto de la clase \textbf{Parser}  un programa Qi inválido, y el parser este comprobará la sintaxis del programa acorde a la arquitectura Qi seleccionada. Al terminar de chequear el programa, el objeto parser retorna un objeto de la clase Exception del tipo \textbf{SyntaxErrorException}, se captura la excepción y se toma de ella el mensaje de error para corroborar que sea el mensaje esperado (que describe la línea donde esta el error de sintáxis).

\end{enumerate}
\subsection{Ensamblado}

Para poder verificar que el ensamblado se realiza correctamente se toma un programa Qi y teniendo en cuenta que el resultado del parser genera un objeto de la clase \textbf{Programa} que contiene un conjunto de objetos de la clase \textbf{Instruccion}, además que cada instrucción tiene el comportamiento necesario para generar su propio \codmaq, se delega la generación del \codmaq a cada una, con la intención de comparar la secuencia resultante con la secuencia de \codmaq esperada.

\subsection{Decodificación}

El proceso de decodificación se verifica a la inversa del proceso de ensamblado: tomando como entrada una lista de cadenas que representa cada una al \codmaq\ de cada instrucción de un programa Qi. Esta lista se recorre para ser procesada por el objeto Interprete que es el encargado de Decodificar. Como resultado de este procesamiento se puede retornar una de dos cosas:
 
\begin{enumerate}
\item \textbf{Lista de Instrucciones} 

Este es el resultado de interpretar una secuencia de código máquina que representa correctamente a un conjunto de instrucciones y sus respectivos operandos. El objeto \textbf{Interprete} verifica los códigos de operación para poder crear las instrucciones correctamente, y luego, dependiendo de que tipo de instrucción se trate, comienza a interpretar sus respectivos operandos (si los tiene).
Por último se compara cada instrucción en la lista resultante con la lista que se esperaba.

\item\textbf{Error}

Al interpretar una secuencia de código máquina que no representa a ninguna instrucción (ya sea porque el código de operación o el código de algún modo de direccionamiento dentro del formato de cada instrucción es invalido) se obtiene una excepción del tipo \textbf{CodigoInvalidoException}.   
\end{enumerate}

\subsection{Ejecución}

Para verificar que la ejecución de un programa Qi se realiza correctamente fue necesario separar la verificación en los tres pasos del \ciclo:

\begin{itemize}
\item \textbf{Búsqueda} 

Para verificar la búsqueda de instrucción se toma la instrucción siguiente a ejecutar y se compara el valor que se guarda en el registro \IR con el valor esperado que debería tener (es decir, con el código maquina que representa a la instrucción siguiente a ejecutar). Ademas se verifica que el registro \PC\ tenga como valor la dirección de memoria de la siguiente instrucción a ejecutar, para corroborar esto se compara el \PC\ luego de la búsqueda con el resultado de la suma del \PC anterior y la cantidad de celdas que ocupa la instrucción procesada.

\item \textbf{Decodificación}

Para verificar la decodificación, se toma el valor que esta guardado registro \IR\ para que lo reciba el interprete que retornará el objeto instrucción interpretado a partir de ese valor. Al obtener dicha instrucción se compara con la instrucción esperada.
 
\item \textbf{Ejecución}

Para verificar el efecto de cada instrucción se preparó un caso de prueba por cada instrucción donde se tiene un estado inicial y un estado final esperado. Se realiza la ejecución de dicha instrucción y se compara el resultado obtenido (ya sea la modificación de flags, registros, celdas de memoria o puertos) con el resultado esperado. 

\end{itemize}

Las siguiente secciones son pasos que se realizan de acuerdo al efecto esperado de cada instrucción en el repertorio de la arquitectura \Q\ (detallada en el apéndice \ref{apendiceQ}).
  
\subsubsection{Operaciones de ALU}

La ALU es la encargada de ejecutar operaciones aritméticas y lógicas. Si la instrucción a ejecutar tiene un efecto aritmético/lógico se le delega la ejecución de la operación. Para poder verificar todas las operaciones creamos un caso de prueba por cada operación. Tanto las operaciones aritméticas como las lógicas se analizan de la misma forma:

\begin{enumerate}
\item Se toma cada valor de los operandos (la búsqueda se detalla el la siguiente sección)

\item Se realiza la operación aritmética/lógica. 

\item Se verifica el resultado obtenido con el resultado esperado.

\item Dado que algunas la mayoría de las operaciones modifica los flags, el valor final de éstos también se comparan con un resultado esperado. 
\end{enumerate}

\subsubsection{Búsqueda de Operandos}

Para verificar la búsqueda de operandos, se toma una instrucción cualquiera y se prueba todas las combinaciones de modos de direccionamiento para cada operando, teniendo en cuenta que son inválidas todas las combinaciones que tienen en el operando destino el modo de direccionamiento \textbf{Inmediato}. Se ejecuta entonces cada combinación para obtener el valor de cada operando, y este valor se lo compara con el valor esperado. 
 
\subsubsection{Almacenamiento de Operandos}

Para verificar el almacenamiento de operandos, se ejecuta cada una de las instrucciones de un programa Qi. Durante la etapa de ejecución, cada una tiene un operando destinado a guardar el resultado de su efecto, y luego de ese momento se puede comparar el valor final del operando con el valor esperado.

%%%%%%%%%%%%%%%%%%%%%%%%%%%%%%%%%%%%%%%%%%%%%%%
\section{Ejemplos de uso}\label{erroresconceptuales}

Durante el desarrollo del simulador se encontraron situaciones especiales de uso que mostraron la potencialidad didáctica de la herramienta, según fue confirmado con los docentes de la materia \orga, dado que los estudiantes pueden experimentar mas allá de la resolución de problemas mediante programas \qi. En esta sección se muestran algunos ejemplos de estas situaciones.

\begin{enumerate}

\item En el siguiente programa \qi\ se ve que el uso de etiquetas permite enmascarar el modo de direccionamiento inmediato difiriendo la resolución de su valor al momento del ensamblado, ya que a partir de entonces es manejado como un valor constante, y por lo tanto puede asignarse a cualquier operando destino. En este ejemplo, la dirección inicial de la rutina \code{rutinaA} se almacena en la celda de memoria \code{0005}, para luego ser recuperada en el \code{CALL}

\begin{verbatim}
rutinaA:  ADD R0, 0x0002 
          ...
          MOV [0x0005], rutinaA
          CALL [0x0005] 
\end{verbatim}

\item En el siguiente ejemplo se puede ejercitar como el uso de subrutinas, puede alterar el \ciclo. En particular en el siguiente ejemplo, como efecto de la ejecución del \code{CALL} se apila la dirección de la siguiente instrucción (en este caso un \code{ADD}), y la ejecución del \code{RET} hace que se vuelva entonces a ejecutar esa misma instrucción y luego un \code{RET} que no tuvo su correspondiente \code{CALL}.

\begin{verbatim}
        CALL rutina 
rutina: ADD R0, 0x0002 
        RET 
\end{verbatim}

%Vamos a hacer un mapa del estado del registro SP antes de seguir: Antes de ejecutar el RET  el registro SP tiene el valor inicial que es FFEF, luego en la ejecución lo primero que se hace es incrementar el SP osea que ahora tiene el valor FFFO y luego buscar el valor de esa dirección para actualizar el registro PC. Para informarles la dirección FFFO es un puerto de E/S. La conclusión es que en esta instancia el flujo de ejecución del programa depende del valor que tiene ese puerto, puede pasar que el valor sea 0x0000 y se actualice el PC nos lleve al inicio de la memoria donde se encuentra inicializado otro programa y empiece a ejecutar desde allí. 


\item Dado el siguiente programa Qi:

\begin{verbatim}
ADD R0, [0x0002]
MUL R4, 0x0001
SUB [0x0003], 0x000A 
\end{verbatim}

Si el valor de PC hace referencia a la ultima instrucción (en este caso un \code{SUB}) y se lleva a cabo el \ciclo, como estado final se tiene el valor de PC apuntando a la siguiente instrucción a ejecutar. Como el programa no tiene mas instrucciones, si se continúa la simulación se comienza la búsqueda de una instrucción inválida.

\end{enumerate} 



%%%%%%%%%%%%%%%%%%%%%%%%%%%%%%%%%%%%%%%%%%%%%%%%%%%
\section{Trabajo Futuro}

En esta sección se describirán características y funcionalidades que deseamos agregar al Simulador QSim en el futuro.

\begin{itemize}

\item Habilitar las instrucciones PUSH y POP.\\

Estas instrucciones permiten el manejo de la Pila como estructura de datos disponible para el programador. La instrucción \code{PUSH} tiene como efecto agregar el valor del operando origen a la pila, mientras que el \code{POP} permite sacar el primer elemento de la pila y guardarlo en el operando destino. Actualmente estas instrucciones se encuentran implementadas pero no habilitadas en ninguna gramática Qi.

\item Entrada y Salida\\

Es deseable que en el futuro el simulador admita la interacción con dispositivos de Entrada/Salida o con simuladores de estos últimos. Para esto habría que modelar buses de entrada hacia los puertos y dispositivos tales como el teclado, impresora, monitor que estén conectados a los puertos y un programa pueda detectar su cambios y utilizar sus datos como input para luego procesarlos en un programa.

\item Implementar interrupciones\\

Para completar el módulo de entrada y salida sería necesario implementar también y agregar al modelo actual las interrupciones de entrada y salida, ya que de esta manera los dispositivos de entrada y salida pueden dar aviso a la cpu que hay nuevos datos para ser procesados.

\end{itemize}


\appendix
\parte{Especificación de la arquitectura Q}\label{apendiceQ}

%%%%%%%%%%%%%%%%%%%%%%%%%%%%%%%%%%%%%%%%%%%%%%%%%%%%%%%%%%%
\section{Características generales} 

La Arquitectura Q tiene 8 registros de uso general de 16 bits, denominados R0..R7, registros especiales de 16 bits tales como PC - \textit{Program counter}, SP\footnote{Comienza en la dirección FFEF.} - \textit{Stack Pointer} y los Flags de un bit como Negative, oVerflow, Carry, Zero. 

Por ultimo tiene una Memoria Principal con direcciones de 16 bit, donde el tamaño de cada celda también es de 16 bit. Entonces la memoria tiene un tamaño de 65536 celdas. \\ 
%%%%%%%%%%%%%%%%%%%%%%%%%%%%%%%%%%%%%%%%%%%%%%%%%%%%%%%%%%%
\section{Modos de direccionamiento}
Los siguientes son los modos de direccionamiento implementados en la Arquitectura \Q.
\begin{enumerate}

\item \textbf{Inmediato} Representa un operando que denota un valor constante. Es importante notar que este modo direccionamiento es admitido en el operando origen pero no el operando destino. La codificación de este modo se indica en la tabla \ref{tablamodos}.

Ejemplos:
\begin{itemize}
\item \textbf{0x0000} denota el modo de direccionamiento inmediato cuyo valor es cero.
\item \textbf{0x000F} denota el modo de direccionamiento inmediato cuyo valor es 15.
\end{itemize}


\item \textbf{Directo}
Con este modo de direccionamiento se denota un operando alojado en una dirección de memoria o de puertos. La codificación de este modo se indica en la tabla \ref{tablamodos}.

Ejemplos:
\begin{itemize}
\item \textbf{[0x0000]} denota un operando cuyo valor se encuentra en la celda de memoria cuya dirección es \textbf{0x0000}.
\item \textbf{[0x000F]} denota  un operando cuyo valor se encuentra en la celda de memoria cuya dirección es \textbf{0x000F}.
\end{itemize}


\item \textbf{Indirecto}
Con este modo de direccionamiento se denota un operando alojado en una celda de memoria cuya dirección está almacenada en otra celda de memoria. La codificación de este modo se indica en la tabla \ref{tablamodos}.

Ejemplos:
\begin{itemize}
\item \textbf{[[0x0000]]} denota un operando cuyo valor se encuentra en la celda de memoria cuya dirección esta guardada como dato en la celda de memoria cuya dirección es \textbf{0x0000}
\item \textbf{[[0x000F]]} denota un operando cuyo valor se encuentra en la celda de memoria cuya dirección esta guardada como dato en la celda de memoria cuya dirección es \textbf{0x000F}
\end{itemize}

\item \textbf{Registro} Con este modo de direccionamiento se denota un operando alojado en un registro de uso general (R0 a R7). La codificación de este modo se indica en la tabla \ref{tablamodos}.

Ejemplos:
\textbf{R0} denota un operando almacenado en el registro R0. 
\textbf{R7} denota un operando almacenado en el registro R7.

\item \textbf{Registro Indirecto} De manera similar al modo indirecto, con este modo de direccionamiento se denota un operando alojado en una celda de memoria cuya dirección está almacenada en el registro indicado. La codificación de este modo se indica en la tabla \ref{tablamodos}.


Ejemplos:
\textbf{[R0]} denota un operando almacenado en una celda de memoria cuya dirección está en el registro \textbf{R0}.
\textbf{[R7]} denota un operando almacenado en una celda de memoria cuya dirección está en el registro \textbf{R7}.

\end{enumerate}

\tabla{|l|l|l|}{
\textbf{Modo}       &  \textbf{Codificación}\\ \hline\hline
Inmediato   &  000000\\ \hline
Directo     &  001000\\ \hline
Indirecto   &  011000\\ \hline
Registro    &  100rrr\\ \hline
Registro indirecto &  110rrr\\ \hline
}{\label{tablamodos}Tabla de códigos de los modos de direccionamiento (Nota: rrr describe el número de registro)}

%%%%%%%%%%%%%%%%%%%%%%%%%%%%%%%%%%%%%%%%%%%%%%%%%%%%%%%%%%%
\section{Repertorio de instrucciones}
En esta sección se detalla cómo se construye el código máquina de las instrucciones de la arquitectura.\\

Todas las instrucciones alteran los flags excepto MOV, CALL, RET, JMP y los saltos condicionales. De las instrucciones que alteran los Flags, todas dejan C y V en 0 a excepción de ADD, SUB y CMP.  \\


%%%%%%%%%%%%%%%%%%%%%%%%%%%%%%%%%%%%%%%%%%%%%%%%%%%%%%%%%%%
\subsection{Instrucciones de 2 operandos}

A continuación se muestra la codificación (formato) de las instrucciones de dos operandos:

\formatoinstr{|c|c|c|c|c|c|}{
Codigo de Operación  & \mdest{}& \msrc{}  & \dest{} & \src{}\\
(4b)     &   (6b)  &  (6b)   &  (16b)  &  (16b)} 


Las instrucciones de dos operandos descriptas a continuación son instrucciones aritméticas o lógicas donde se asume que el resultado de la operación se almacena en uno de los dos operandos de entrada, y por lo tanto se lo denomina \textbf{operando destino}.

\begin{enumerate}
\item  \mostrarinst{MUL destino, origen}{0000}\\
Esta instrucción describe la multiplicación entre los datos de los dos operandos. Esta operación es la única que cuyo resultado puede ser 32 bits, que son lo que ocuparía más de una celda de memoria en código binario, por lo que los primeros 16 bits, es decir, la primer mitad, es guardada en el registro \textbf{R7} y la segunda en el operando destino.
 
\item  \mostrarinst{ADD destino, origen}{0010}\\
Esta instrucción describe la suma entre los datos de los dos operandos. El resultado de la ejecución de la suma es guardado en el operando destino.

\item  \mostrarinst{SUB destino, origen}{0011}\\
Esta instrucción describe la resta entre los datos de los dos operandos. El resultado de la ejecución de dicha resta es guardado en el operando destino.

\item  \mostrarinst{DIV destino, origen}{0111}\\
Esta instrucción describe la división entre el dato en el operando destino como dividendo y el dato en el operando origen como divisor. El resultado de la ejecución de la división es guardado en el operando destino.

\item  \mostrarinst{MOV destino, origen}{0001}\\
Esta instrucción describe la copia de datos del dato alojado en el operando origen al operando destino. El resultado de la ejecución del MOV es el dato guardado en el operando origen ahora también guardado en el operando destino.

\item  \mostrarinst{AND destino, origen}{0100}\\
Esta instrucción describe la conjunción bit a bit entre los datos de los dos operandos. El resultado de la ejecución de esta operación es guardado en el operando destino.

\item  \mostrarinst{CMP destino, origen}{0110}\\
Esta instrucción describe la resta entre dos operandos, sin guardar el resultado. Su único efecto es la actualización de flags en la cpu.

\item  \mostrarinst{OR destino, origen}{0101}\\
Esta instrucción describe la disyunción bit a bit entre los datos de los dos operandos. El resultado de la ejecución de esta operación es guardado en el operando destino.
\end{enumerate}

%%%%%%%%%%%%%%%%%%%%%%%%%%%%%%%%%%%%%%%%%%%%%%%%%%%%%%%%%%%
\subsection{Instrucciones de 1 operando origen}

El formato de las instrucciones de un operando origen es el siguiente:\\

\formatoinstr{|c|c|c|c|}
{CodOp& Relleno&\msrc{}& Operando \src{}\\
(4b)     &   (000000)  &  (6b)   &  (16b) 
}

\begin{enumerate}
\item  \mostrarinst{CALL origen}{1011}\\
El efecto del CALL es guardar la dirección de memoria en la celda de la dirección que se encuentra guardada en el SP (Stack pointer) aumentar el SP y guardar en el PC (Program Counter) el dato que se encuentra guardado en el operando origen ya que describe el llamado a una subrutina que comienza en la celda de memoria cuya dirección esta guardada en el operando origen.

\item  \mostrarinst{JMP origen}{0110}\\
El efecto del JMP es cambiar el PC (Program Counter) por el dato que esta guardado en el operando origen ya que esta operación describe el salto a otra parte de la memoria para continuar con la ejecución del programa.
\end{enumerate}

%%%%%%%%%%%%%%%%%%%%%%%%%%%%%%%%%%%%%%%%%%%%%%%%%%%%%%%%%%%
\subsection{Instrucciones de 1 operando destino}

El formato de las instrucciones de un operando destino es el siguiente:\\

\formatoinstr{|c|c|c|c|}
{CodOp & \msrc{} & Relleno & Operando \src{}\\
(4b)   &   (6b)  & (000000)&  (16b) 
} 

\begin{enumerate}
\item  \mostrarinst{NOT destino}{1001}\\
Esta instrucción describe la operación lógica "negación" bit a bit en el datos del operando destino. El resultado de la ejecución de esta operación es guardado en la misma celda o registro de donde es leído el dato inicialmente.
\end{enumerate}

%%%%%%%%%%%%%%%%%%%%%%%%%%%%%%%%%%%%%%%%%%%%%%%%%%%%%%%%%%%
\subsection{Instrucciones sin operandos}

El formato de las instrucciones sin operandos es el siguiente:\\

\formatoinstr{|c|c|}{CodOp & Relleno \\  (4b) & (000000000000)}


\begin{enumerate}
\item  \mostrarinst{RET}{0110}
El efecto del ret es cambiar el pc por el dato que esta guardado en la celda de memoria que se encuentra en el SP (Stack pointer) y decrementar el SP ya que describe la finalización de la ejecución de una subrutina y la ejecución del resto del programa.
\end{enumerate}


%%%%%%%%%%%%%%%%%%%%%%%%%%%%%%%%%%%%%%%%%%%%%%%%%%%%%%%%%%%
\subsection{Instrucciones de salto condicional  (falta revisar Mara)}

El formato de las instrucciones de salto condicional es el siguiente:\\

\formatoinstr{|c|c|}{CodOp (8)&Desplazamiento(8)} donde los primeros cuatro bits del campo \code{CodOp} es la cadena \code{1111}. Si \textbf{al evaluar la condición de salto} el resultado es 1, se le suma al PC el valor del desplazamiento, representado en el sistema Complemento a 2 de 8 bits. En caso contrario la instrucción no hace nada.


El efecto de cualquier salto condicional es aumentar el PC (Program Counter) en la cantidad de celdas que indique el desplazamiento si sólo la condición que cada salto condicional tiene da como resultado 1, lo cual es interpretado como verdadero.

\begin{enumerate}
\item  \mostrarinst{JE desplazamiento}{0001}\\
La condición del salto es que el flag \textbf{Z} (Cero) sea 1, es decir la ultima operación matemática dió como resultado el número cero.

\item  \mostrarinst{JNE desplazamiento}{1001}\\
La condición del salto es que el flag \textbf{Z} (Cero) sea 0, es decir la ultima operación matemática no dió como resultado el número cero.

\item  \mostrarinst{JLE desplazamiento}{0010}\\
La condición del salto es el resultado de la siguiente operación lógica \textbf{Z OR ( N XOR V )}, es decir la ultima operación matemática es menor o igual con signo.

\item  \mostrarinst{JG desplazamiento}{1010}\\
La condición del salto es el resultado de la siguiente operación lógica \textbf{NOT (Z OR ( N XOR V ))}, es decir la ultima operación matemática es mayor con signo.

\item  \mostrarinst{JL desplazamiento}{0011}\\
La condición del salto es el resultado de la siguiente operación lógica \textbf{N XOR V}, es decir la ultima operación matemática es menor con signo.

\item  \mostrarinst{JGE desplazamiento}{1011}\\
La condición del salto es el resultado de la siguiente operación lógica \textbf{NOT (N XOR V)}, es decir la ultima operación matemática es mayor o igual con signo.

\item  \mostrarinst{JLEU desplazamiento}{0100}\\
La condición del salto es el resultado de la siguiente operación lógica \textbf{C OR Z}, es decir la ultima operación matemática es menor o igual sin signo.

\item  \mostrarinst{JGU desplazamiento}{1100}\\
La condición del salto es el resultado de la siguiente operación lógica \textbf{NOT (C OR Z)}, es decir la ultima operación matemática es mayor sin signo.

\item  \mostrarinst{JCS desplazamiento}{0101}\\
La condición del salto es que el flag \textbf{C} sea 1, es decir la ultima operación matemática es menor sin signo.

\item  \mostrarinst{JNEG desplazamiento}{0101}\\
La condición del salto es que el flag \textbf{N} sea 1, es decir si el último resultado de una operación dio negativo.

\item  \mostrarinst{JVS desplazamiento}{0111}\\
La condición del salto es que el flag \textbf{V} sea 1, es decir si el último resultado de una operación dio overflow.

\end{enumerate}
\section{Como utilizar el simulador}

En esta sección se describen las opcciones básicas del uso de la aplicación.

\subsubsection{Agregar archivos}
Como se observa en screenshot la ventana de cargado la misma cuenta con la opcción de agregar archivos .qsim y se encuentran desabilitadas las otras opcciones ya que, sin uno o mças programas de la arquitectura Qi carece de sentido realizar la acción de ensamblar o de cargar en memoria.
\graf{ventana_cargado}{Cargado de archivos}

\subsubsection{Ensamblar}
Como se observa en screenshot de la ventana Ensamblar una vez que los programas Qi en los archivos -qsim se encuentran agregados se habilitan las opcciones de seleccionar la arquitectura Qi que se desee (Q1.. Q6) y la opcción de Ensamblar el programa para que se realice el chequeo de sintáxis y se genere el código máquina que luego será cargado en memoria.

\graf{ventana_cargado_ensambar}{Ensamblar}

\subsubsection{Cargar en memoria}
Como se observa en screenshot Cargado en memoria los programas Qi en los archivos -qsim se encuentran ensamblados en la arquitectura Qi elegida.
Para ese entonces todas las opcciones estan habilitadas ya que puede desearse agregar otro archivo .qsim o quitar alguno (no todos) y volver a ensamblar, o bien, elegir el \PC a partir del cual se quiere cargar en memoria el programa Qi (por defecto "0000"') y hacer click en el botón "Cargar en memoria".
\graf{ventana_cargado_cargarenmemoria}{Cargado en memoria}

\subsubsection{Ciclo de instruccion}
Como se observa en screenshot de la ventana de ejecución una vez que el programa es exitosamente cargado en memoria se abre esta nueva ventana que contiene los registro especiales \PC, \IR, \SP, los flags, los registro de uso general (R0...R7), la memoria, una consola de información al usuario y los botones habilitados para realizar el fetch, ver los puertos y editar los valores de los registros especiales, los flags y pc.
\graf{ventana_ejecucion}{Ventana de ejecución}\\
Para poder realizar el ciclo de ejecución de una instruccion, se debe seguir el orden de las etapas de dicho ciclo al apretar los botones que estan debajo en el panel llamado "Ciclo de ejecución". De todas maneras como se ve en los siguientes screenshots sólo luego de realizarse el fecth a pedido del usario puede utilizarse el decode, luego el execute y el ciclo vuelve a repetise.\\

\graf{ventana_ejecucion_fetch}{Ventana de ejecución luego del fetch}\\
\graf{ventana_ejecucion_decode}{Ventana de ejecución luego del decode}\\
\graf{ventana_ejecucion_execute}{Ventana de ejecución luego del execute}\\

\subsubsection{Ver puertos}
Si se desea ver el valor que tienen los puertos, basta con hacer click en el botón "Ver puertos", y se abrirá una nueva ventana tal y como se ve en el screenshot de la ventana de puertos donde se podrá observar el número de puerto y su valor actual y también al igual que en la ventana principal seleccionar su edición.
\graf{ventana_puertos}{Ventana de puertos}\\


\section{Repositorios remotos del código de la aplicación}

\subsubsection{Código del modelo}\\
https://github.com/molinarirosito/QSim

\subsubsection{Código de la interfaz}\\
https://github.com/molinarirosito/QSim_UI
\section{Errores comunes}


En esta sección se detallan las diferentes situaciones que pueden dar como resultado un mensaje de error, que informará en que línea del programa se encuentra, durante la etapa de ensamblado.

\begin{itemize}

\item Dado el siguiente programa Qi:

\begin{verbatim}
ADD R0, 0x0002
MUL R4, 0x01
SUB R5, 0x000A
MOV R5, 0x0056
MOV R2, R3
ADD R1, R7
\end{verbatim}

En la linea numero 2 el modo de direccionamiento inmediato esta incompleto, (le faltan dos dígitos). Cuando se quiera ensamblar este programa, el ensamblador detectará este error y en la pantalla se mostrará mediante un pop-up\footnote{Ventana emergente.} el siguiente mensaje:
 
\mensaje{Ha ocurrido un error en la linea 2 : MUL R4, 0x01} 

\item Dado el siguiente programa Qi:

\begin{verbatim}
MOV 0x0006, 0x0056
ADD R2, R3
SUB R1, R7
\end{verbatim}

En la linea numero 1 el operando destino es inmediato lo cual es invalido (El operando nunca puede tener como modo de direccionamiento un inmediato). Cuando se quiera ensamblar este programa, el ensamblador detectará este error y en la pantalla se mostrará mediante un pop-up el siguiente mensaje:

\mensaje{Ha ocurrido un error en la linea 1 : MOV 0x0006, 0x0056} 

\item Dado el siguiente programa escrito en Q1:

\begin{verbatim}
ADD R0, [0x0002]
MOV R4, R0
\end{verbatim}

En la linea numero 1 el operando origen es directo lo cual es invalido en la arquitectura Q1 (El modo de direccionamiento Directo se incorpora en las arquitecturas Qi desde la arquitectura Q2 en adelante). Cuando el alumno quiera ensamblar este programa, el ensamblador detectará este error y en la pantalla se mostrará mediante un pop-up el siguiente mensaje:

\mensaje{Ha ocurrido un error en la linea 1 : ADD R0, [0x0002]} 

\item Dado el siguiente programa Qi:

\begin{verbatim}
CMP R3, [0xA000]
MOV R4 R0 
\end{verbatim}

En la línea numero 2 entre los operandos no se encuentra la coma que los separa (La sintáxis de las instrucciones de dos operandos especifica que debe haber una coma separando los operandos.). Cuando se quiera ensamblar este programa, el ensamblador detectará este error y en la pantalla se mostrará mediante un pop-up el siguiente mensaje:

\mensaje{Ha ocurrido un error en la linea 2 : MOV R4, R0} 

\item Dado el siguiente programa Qi:

\begin{verbatim}
sub [[0x0004]], [0xA000]
ADD R4, R0
\end{verbatim}

En la linea numero 1 la instrucción sub esta escrita en minúscula esto es inválido (La sintáxis define que los nombres de las instrucciones son estrictamente en mayúscula). Cuando se quiera ensamblar este programa, el ensamblador detectará este error y en la pantalla se mostrará mediante un pop-up el siguiente mensaje:

\mensaje{Ha ocurrido un error en la linea 1: sub [[0x0004]], [0xA000]}

\item Dado el siguiente programa Qi:

\begin{verbatim}
MUL [R6], r4
ADD [0xF0F0], R0
\end{verbatim}

En la linea numero 1 el operando origen es un registro escrito con minúscula, esto es invá lido(La sintáxis define que los registros empiezan estrictamente con 'R'' mayúscula). Cuando se quiera ensamblar este programa, el ensamblador detectara este error y en la pantalla se mostrará mediante un pop-up el siguiente mensaje:

\mensaje{Ha ocurrido un error en la linea 1: MUL [R6], r4} 

\item Dado el siguiente programa Qi:

\begin{verbatim}
AND R2, R8
OR [0xF0F0], R0
\end{verbatim}

En la linea numero 1 el operando origen el numero del registro es invalido (Los registros deben estar dentro del rango R0..R7). Cuando se quiera ensamblar este programa, el ensamblador detectará este error y en la pantalla se mostrará mediante un pop-up el siguiente mensaje:

\mensaje{Ha ocurrido un error en la linea 1: AND R2, R8} 

\item Dado el siguiente programa Qi:

\begin{verbatim}
MUL R7, R4
AND R5, [R3]
\end{verbatim}

Para la instrucción MUL es inválido utilizar como destino el registro R7 por lo que, la primer línea es inválida. Cuando se quiera ensamblar este programa, el ensamblador detectará este error y en la pantalla se mostrará mediante un pop-up el siguiente mensaje:

\mensaje{Ha ocurrido un error en la linea 1: MUL R7, R4}  

\item Dado el siguiente programa Qi:

\begin{verbatim}
inicio: MUL R7, R4
        AND R5, [R3]
        JMP incio
\end{verbatim}

En la linea numero 3 la etiqueta anteriormente declarada en la línea número 1 esta incompleta (Le faltan los dos puntos). Cuando se quiera ensamblar este programa, el ensamblador detectará este error y en la pantalla se mostrará mediante un pop-up el siguiente mensaje:

\mensaje{Ha ocurrido un error en la linea 3: JMP incio}


\item Dado el siguiente programa Qi :

\begin{verbatim}
ADD [0x9000], R4
NOT 0x0004
\end{verbatim}

En la linea numero 2 el operando destino de la instrucción NOT no puede ser un inmediato, esto es inválido (Los operandos destinos no pueden ser inmediatos). Cuando se quiera ensamblar este programa, el ensamblador detectará este error y en la pantalla se mostrará mediante un pop-up el siguiente mensaje:

\mensaje{Ha ocurrido un error en la linea 2: NOT 0x0004}  

\item Dado el siguiente programa Qi:

\begin{verbatim}
ADD [9000], R4
NOT R2
\end{verbatim}

En la linea numero 1 el operando origen no tiene el prefijo '0x', es una expresión invalida. Cuando se quiera ensamblar este programa, el ensamblador detectará este error y en la pantalla se mostrará mediante un pop-up el siguiente mensaje:

\mensaje{Ha ocurrido un error en la linea 1: ADD [0009], R4} 

\item Dado el siguiente programa Qi:

\begin{verbatim}
ADD [0x900000000000], R4
NOT R2
\end{verbatim}

En la linea numero 1 el operando destino tiene mas dígitos que los permitidos, (Un inmediato tiene el prefijo 0x y luego sólo 4 dígitos hexadecimales). Cuando se quiera ensamblar este programa, el ensamblador detectará este error y en la pantalla se mostrará mediante un pop-up el siguiente mensaje:

\mensaje{Ha ocurrido un error en la linea 1: ADD [0x900000000000], R4}  

\item Dado el siguiente programa Qi:

\begin{verbatim}
ZDD [0x900000000000], [R5]
NOT R2
\end{verbatim}

En la linea numero 1 el nombre de la operación es invalida (No existe la instrucción ZDD).Cuando se quiera ensamblar este programa, el ensamblador detectará este error y en la pantalla se mostrará mediante un pop-up el siguiente mensaje:

\mensaje{Ha ocurrido un error en la linea 1: ZDD [0x900000000000], R4}

\item Dado el siguiente programa Qi:

\begin{verbatim}
SUB [], 0x000A
\end{verbatim}

En la linea numero 1 el operando destino no esta incompleto (El modo de direccionamiento directo debe escribirse como un valor inmediato encerrado entre corchetes). Cuando se quiera ensamblar este programa, el ensamblador detectará este error y en la pantalla se mostrará mediante un pop-up el siguiente mensaje:

\mensaje{Ha ocurrido un error en la linea 1: SUB [], 0x000A} 

\item Dado el siguiente programa Qi:

\begin{verbatim}
JMP
\end{verbatim}

En la linea numero 1 sólo esta escrito el nombre de la instruccion JMP pero falta a continuación su el operando origen. Cuando se quiera ensamblar este programa, el ensamblador detectará este error y en la pantalla se mostrará mediante un pop-up el siguiente mensaje:

\mensaje{Ha ocurrido un error en la linea 1: JMP} 
 

\end{itemize}

\newpage
%%%%%%%%%%%%%%%%%%%%%%%%%%%%%%%%%%%%%%%%%%%%%%%%%%%%
\begin{thebibliography}{9}

\bibitem{Stallings} Williams Stallings, \emph{Computer Organization and Architecture}, octava edición, Editorial Prentice Hall, 2010.

\bibitem{Tanenbaum} A. Tanenbaum, \emph{Organización de Computadoras}, cuarta edición, Editorial Pearson.

\bibitem{Patterson} Hennessy, Patterson. \emph{Arquitectura de Computadores - Un enfoque cuantitativo}, primera edición,  Editorial Mc Graw Hill.

\bibitem{blog} Sitio oficial de la materia \orga: \url{http:\\orga.blog.unq.edu.ar} (2013)

\end{thebibliography}

\end{document}